\section{Preparación, pilotos y habilitación de innovaciones}
Los pilotos conservan la correspondencia entre paquetes 1.9 del T-14, actividades y cargas del T-15 y códigos IN de SD13. Se ejecutarán con usuarios autorizados, versiones identificadas, responsables de soporte y ventana de intervención permitida. La aceptación de una innovación no sustituye el cierre de marcha blanca de la solución obligatoria ni sus pruebas de volumen, disponibilidad, respuesta, capacitación y conciliación del artículo 17.3 \parencite[arts.~17--18; pp.~12--13]{bases-admin}.

IN1 levantará su registro base en el mes 18 con el flujo adicional desactivado para la cohorte y lo evaluará en el 19, después de capacitación. La primera respuesta incluye aceptación o rechazo y se conservarán solicitudes abiertas; se comprobará el historial de todas las transiciones. Una mejora sin universo comparable no habilita declarar beneficio. El paso a operación exige decisión y transferencia del paquete 1.9.1.6.1.

IN2 requiere aceptación técnica, IN1 habilitada, participantes, fichas y planes confirmados. Su piloto conserva 31 a 36; la inscripción y los planes se preparan en 30 y la evaluación y transferencia cierran en 37: observación invernal junio--agosto y botaduras septiembre--noviembre. Los informes se reciben por quincena; la revisión de cada embarcación se programa tres semanas antes de su botadura y se dimensiona por los cinco lotes 1.9.2.5.4.k. Las restricciones de campaña se aplican a cualquier intervención física; el seguimiento documental no autoriza instalación ni trabajos de varadero dentro de una ventana bloqueada.

IN3 requiere muestra autorizada, controles de IA y usuarios preparados antes de comparar ejecución manual y asistida en 18--19. La revisión humana confirmará los campos antes del registro definitivo; el ingreso manual seguirá disponible ante error o indisponibilidad del extractor. Los casos rechazados, ilegibles y corregidos forman parte de la evaluación. La habilitación exige informe de tiempos, errores, costo y decisión, y transferencia 1.9.3.6.1.

IN4 conservará captura propia del tiempo activo y comparación de 18--19 coordinada con IN1, sin duplicar evidencia ni beneficio. El piloto no genera pagos variables. La regla de cálculo, calidad y atribución deberá estar aceptada antes de reconocer el incentivo en operación; ausencia de evidencia admisible produce factor cero. El expediente contractual se prepara y aprueba con funciones separadas.

IN5 requiere medición individual validada, permanencias conciliadas, consentimiento vigente y adhesión voluntaria. La comparación usa 28 días de base en 18 y 28 de evaluación en 19; la acción y sus evidencias se sitúan entre esas ventanas. Las vistas mostrarán cobertura y último corte válido, y no atribuirán una reducción a intervalos incompletos o incomparables. El seguimiento continuará con el informe mensual de 1.12.2.5.m y evidencia anterior a la temporada 2029, respetando las habilitaciones del SD2.

Cada habilitación registrará versión, cohorte, permisos, datos de entrada, indicadores, incidencias, decisión y responsable. Ante acceso indebido, diferencias de datos o pérdida de trazabilidad se detendrá la ampliación, se desactivará el flujo adicional y se conservarán solicitudes y decisiones ya emitidas. La recuperación reconciliará esas escrituras antes de reabrir el servicio; desactivar una función no equivale a borrar registros aceptados. La matriz de reversión distinguirá datos base e incrementales y se comprobará junto con los procedimientos de SD5. La reanudación requerirá corrección comprobada y autorización, sin eliminar los controles de recepción.

